\section{Notation and Minimal Assumptions}


The symbols used in this paper constitute a local working notation; they do not carry over the standard meanings of existing theories.

\begin{itemize}
\item $t$: physical time, externally measurable.
\item $\tau$: internal time delay, buffering, retention. Not subjective time itself.
\item $v_f$: fast, predictive and reactive dynamics.
\item $v_s$: slow, sedimentary and stabilizing dynamics.
\item $\Delta v = v_f - v_s$: speed difference.
\item $\mu$: coupling viscosity between speed scales.
\item $\mathcal{T}$: effective torque-like quantity that transforms speed difference.
\item $\Phi$: information potential.
\item $I$: information density.
\item $J_{\log}$: log-referencing flow in the consciousness window.
\end{itemize}

Crucially, $\tau$ is not used as torque. $\tau$ denotes internal delay; the torque-like quantity is consistently written $\mathcal{T}$. Both quantities involve $\Delta v/\mu$ in their structural expressions, but they describe different aspects: $\tau$ is the temporal thickness of referencing, while $\mathcal{T}$ is the transformation load on speed difference.

The structural relations among these quantities are organized in §7.
